\abstract % Структурный элемент: РЕФЕРАТ. Строка со сведениями об объёме формируется автоматически.

\keywords{APACHE PAIMON, КОЛОНОЧНЫЙ ФОРМАТ, VORTEX, PARQUET, ORC, LAKEHOUSE, LSM-ДЕРЕВО, КОДИРОВАНИЕ ДАННЫХ, ПРОПУСК ДАННЫХ, НАГРУЗОЧНОЕ ТЕСТИРОВАНИЕ, APACHE FLINK}

Объект исследования --- колоночные файловые форматы хранения данных в составе
lakehouse-системы Apache Paimon: Apache Parquet, Apache ORC и Vortex. Предмет
исследования --- алгоритмы кодирования, сжатия и пропуска данных, реализуемые
этими форматами, и их влияние на производительность чтения и записи таблиц с
первичным ключом и таблиц только для добавления.

Цель работы --- определить область применимости формата Vortex в Apache Paimon,
разработать модуль его интеграции и провести сравнительный анализ форматов на
представительных нагрузках.

В ходе работы выполнен аналитический обзор эволюции систем хранения аналитических
данных --- от хранилищ данных к архитектуре lakehouse, --- открытых табличных
форматов и колоночных файловых форматов; разработан и интегрирован в Apache Paimon
модуль поддержки формата Vortex: реализация интерфейса \texttt{FileFormat},
конвертер предикатов, взаимодействие с нативной библиотекой через Java Native
Interface; проведено функциональное тестирование артефакта в контейнерах;
развёрнуто окружение нагрузочного тестирования на кластере из пяти узлов
средствами Ansible; проведены две серии экспериментов --- batch-аналитика на
наборе TPC-DS масштаба \SI{100}{\giga\byte} и near-real-time-сценарий с потоковой
записью и параллельными аналитическими запросами; собраны и статистически
обработаны метрики производительности.

Результаты: на batch-аналитике (таблицы только для добавления, крупные файлы)
Vortex доминирует над Parquet --- быстрее на \num{82} запросах из \num{103}, среднее
ускорение порядка \SI{25}{\percent}, пиковое --- до \SI{80}{\percent}; именно этот
сценарий --- основная область применимости Vortex. На потоковой
near-real-time-нагрузке (таблица с первичным ключом, слияние при чтении) Vortex
эквивалентен Parquet и ORC по пропускной способности записи; по типичной задержке
чтения (медиане) Vortex --- быстрейший формат в обеих фазах и на большинстве
запросов (девять из тринадцати как в фазе FINAL, так и в фазе DURING), особенно ---
на запросах с селективным предикатом, диапазонным фильтром по строке-префиксу и в
многошаговой аналитике. На простой группировке и многомерной агрегации без
селективного предиката Vortex незначительно (в пределах единиц процентов) уступает
конкурентам --- следствие отсутствия проталкивания агрегатов в интерфейсе форматов
Apache Paimon. По хвостам задержки в фазе DURING ORC и Vortex показывают узкие
распределения (среднее почти равно медиане), Parquet --- более широкое
распределение с эпизодическими выбросами в десятки секунд в моменты сброса данных
по чекпойнту, что согласуется с архитектурным выигрышем Vortex за счёт независимого
нативного пула соединений с объектным хранилищем. Сформулированы рекомендации по
выбору формата хранения для различных классов нагрузок.

Область применения результатов --- проектирование витрин данных на базе Apache
Paimon, настройка lakehouse-хранилищ, развитие экосистемы Apache Paimon.
